% Aqui começa o capítulo de Análise dos Domínios de Desempenho.
\chapter{Análise dos Domínios de Desempenho}\label{chp:DominiosDesempenho}

O PMBOK 7ª edição define oito domínios de desempenho que representam áreas críticas a serem continuamente observadas e ajustadas ao longo do projeto. Esses domínios não correspondem a etapas sequenciais, mas a dimensões interdependentes que influenciam o sucesso do projeto de forma simultânea. Nesta seção, cada domínio é analisado à luz do projeto de desenvolvimento do sistema de gestão de casos sociais, considerando a adoção do framework Scrum como abordagem principal.

\section{Domínio dos Stakeholders}

O domínio dos stakeholders trata da identificação, do engajamento e da gestão das expectativas das partes interessadas ao longo do projeto. No contexto de um sistema de gestão de casos sociais, os stakeholders incluem equipes técnicas, gestores institucionais, coordenadores de políticas públicas e usuários operacionais do sistema.

A abordagem ágil favorece o engajamento contínuo desses atores por meio de interações frequentes, especialmente durante as revisões de sprint. Esse mecanismo permite alinhar expectativas, validar incrementos do produto e ajustar prioridades de acordo com o feedback recebido. O desempenho nesse domínio é avaliado pela qualidade do engajamento, pelo nível de participação dos stakeholders e pela aderência das entregas às necessidades identificadas.

\section{Domínio da Equipe}

O domínio da equipe enfatiza a importância de equipes capacitadas, colaborativas e auto-organizadas. O desenvolvimento do sistema proposto exige a integração de competências técnicas, compreensão do domínio social e capacidade de adaptação a mudanças frequentes.

No Scrum, a equipe de desenvolvimento atua de forma multifuncional, compartilhando responsabilidades sobre o resultado do sprint. O desempenho nesse domínio está associado à comunicação eficaz, à resolução colaborativa de problemas e à capacidade da equipe de manter um ritmo sustentável de trabalho ao longo do projeto.

\section{Domínio da Abordagem de Desenvolvimento e do Ciclo de Vida}

Este domínio trata da escolha e da adaptação da abordagem de desenvolvimento e do ciclo de vida do projeto. Dada a natureza do sistema de gestão de casos sociais, caracterizada por requisitos emergentes e necessidade de validação constante com usuários, foi adotada uma abordagem ágil baseada no Scrum.

O ciclo de vida incremental e iterativo permite entregar funcionalidades de forma progressiva, reduzindo riscos e incorporando aprendizado contínuo. O desempenho nesse domínio é observado pela adequação da abordagem escolhida ao contexto do projeto e pela capacidade de ajustar práticas conforme o projeto evolui.

\section{Domínio do Planejamento}

No PMBOK 7, o planejamento é entendido como uma atividade contínua, adaptativa e orientada a objetivos, e não como uma etapa fixa realizada apenas no início do projeto. No projeto em análise, o planejamento ocorre em múltiplos níveis, incluindo a definição de uma visão do produto, a elaboração de um roadmap e o refinamento contínuo do backlog.

O Scrum operacionaliza esse domínio por meio de eventos como o planejamento da sprint e o refinamento do Product Backlog. O desempenho nesse domínio é avaliado pela clareza dos objetivos, pela priorização adequada do trabalho e pela capacidade de responder a mudanças sem perda de foco estratégico.

\section{Domínio do Trabalho do Projeto}

O domínio do trabalho do projeto refere-se à execução coordenada das atividades necessárias para produzir os resultados esperados. No desenvolvimento do sistema de gestão de casos sociais, esse domínio envolve a implementação das funcionalidades, a integração dos componentes e a validação dos incrementos entregues.

A transparência promovida pelos artefatos do Scrum, como o Sprint Backlog e os quadros de acompanhamento visual, contribui para o monitoramento contínuo do trabalho. O desempenho nesse domínio está relacionado à fluidez do trabalho, à identificação precoce de impedimentos e à capacidade de manter a consistência das entregas ao longo das sprints.

\section{Domínio da Entrega}

O domínio da entrega concentra-se na disponibilização efetiva de valor aos stakeholders. Em projetos ágeis, a entrega ocorre de forma incremental, permitindo que partes funcionais do produto sejam utilizadas ou avaliadas antes da conclusão total do projeto.

No sistema de gestão de casos sociais, a entrega de incrementos como o cadastro básico de casos, a categorização por tags e a visualização em mapa permite validar o valor da solução desde as fases iniciais. O desempenho nesse domínio é medido pela utilidade dos incrementos entregues, pela satisfação dos stakeholders e pela capacidade de gerar benefícios mensuráveis.

\section{Domínio da Medição}

A medição é essencial para apoiar a tomada de decisão e o aprendizado ao longo do projeto. O PMBOK 7 enfatiza a utilização de métricas relevantes, alinhadas aos objetivos do projeto e ao contexto organizacional.

No projeto proposto, são utilizadas métricas ágeis, como progresso do backlog, previsibilidade das sprints e indicadores de fluxo de trabalho. Essas métricas permitem avaliar o desempenho do projeto, identificar tendências e apoiar ajustes no planejamento e na execução, sem recorrer a controles excessivamente burocráticos.

\section{Domínio da Incerteza}

O domínio da incerteza reconhece que projetos operam em ambientes nos quais nem todas as variáveis podem ser conhecidas antecipadamente. No desenvolvimento de um sistema de gestão de casos sociais, a incerteza está presente nos requisitos, no comportamento dos usuários e em possíveis mudanças institucionais ou regulatórias.

A adoção do Scrum permite tratar a incerteza de forma explícita, por meio de ciclos curtos de aprendizado, validação frequente e adaptação contínua. O desempenho nesse domínio é observado pela capacidade do projeto de absorver mudanças, mitigar riscos e evoluir a solução com base em evidências obtidas ao l

